% !TeX root = main.tex
\section{数据科学方法}

\paragraph{示例：木材分拣问题}

设想一家锯木厂希望用自动化设备将混杂在一起的木板按树种分拣开来，例如区分白蜡木（Esche）与桦木（Birke），如图 \ref{fig:wood} 所示。整套技术方案由三个环节组成：摄像头依次拍下每一块木板的图像，特征提取环节把图像信息压缩为若干具有区分度的数值特征，分类环节再根据这些特征做出判断，最终输出"白蜡木"或"桦木"的决策。用符号表示，这个过程是从图像空间 $\mathbb{R}^m \times \mathbb{R}^n \times \cdots$，经特征提取映射到特征空间 $\mathbb{R}^d$，再经分类映射到类别集合 $\{Ash, Birch\}$。

\begin{figure}[htbp]
  \centering
  \includegraphics[width=0.5\textwidth]{wood.png}
  \caption{白蜡木与桦木}
  \label{fig:wood}
\end{figure}

最朴素的做法，是凭经验假设白蜡木的颜色比桦木更深，于是用图像的平均亮度 $x$ 作为唯一特征，选定一个阈值 $x_S$，亮度超过阈值就判为桦木，反之判为白蜡木；阈值本身由少量样本图像估计得出。用真实样本画出的亮度直方图 \ref{fig:woodhistogram} 确实印证了"白蜡木更暗"这一经验，但两类样本的亮度分布存在明显重叠，单靠一个亮度阈值无法稳健地区分两种木材。

\begin{figure}[htbp]
  \centering
  \includegraphics[width=0.5\textwidth]{woodhistogram.png}
  \caption{两种木头的亮度直方图}
  \label{fig:woodhistogram}
\end{figure}

于是很自然地想到引入第二个特征：用频谱分析衡量木纹的强弱程度，记作 $x_2$，与亮度 $x_1$ 一起构成二维特征向量 $\mathbf{x} = (x_1, x_2)^t$。把带标签的样本画在这个二维特征空间中，如图 \ref{fig:wood2d}，可以看到两类样本此时已经能被一条直线（记作 $AB$，即决策边界）较好地分开。但随之而来的问题是：这条边界对新样本是否依然有效？它又该如何通过训练自动获得，而不是靠人眼观察画出来？

\begin{figure}[htbp]
  \centering
  \includegraphics[width=0.5\textwidth]{wood2d.png}
  \caption{二维特征空间}
  \label{fig:wood2d}
\end{figure}

这个小例子浓缩了模式分类要处理的核心矛盾：选择怎样的特征、采用怎样的模型、以及如何用有限的样本训练出一个能够推广到未见样本的决策规则。

\subsection{总体设计流程}

围绕上述问题，可以归纳出一套通用的设计循环：采集数据、选择特征、选择模型、训练分类器、评估分类器，评估结果又反过来指导对特征、模型乃至数据采集的调整，如此循环，直至性能满足要求。其中选择模型这一步会用到关于问题的先验知识（例如某种不变性），而训练分类器在不同语境下也被称作学习或优化。

要让训练出的分类器具备推广能力，样本集合通常被划分为三部分：训练样本用于直接训练分类器；独立的验证样本用于决定何时停止训练；独立的测试样本用于最终评估分类器的性能。若拿训练集本身去评估分类器，则会得到过于乐观、不可信的结果，这在方法论上被视为一种错误。

训练过程中，把全部训练样本完整地送入一次称为一个周期（epoch）。根据样本呈现方式的不同，训练协议可分为批量训练（在学习开始前一次性给出全部样本）、序贯训练（样本逐个给出）、随机训练（样本被视为随机变量）以及在线训练（每个样本只呈现一次）。

分类误差随训练过程的变化可以用学习曲线来观察：训练集上的平均误差 $E$ 会逐渐逼近某个最小值 $E_{\min}$，这个最小值取决于贝叶斯误差本身、训练数据量以及分类器 $C_p(x)$ 的表达能力。实践中常见的做法是在独立验证集误差出现第一个极小值时就停止训练，因为继续训练下去，验证误差往往会重新上升，这正是过拟合的信号；测试集上的误差平均而言也会高于训练集上的误差。

分类器 $C_p(\mathbf{x}_i)$ 对样本 $\mathbf{x}_i$ 给出类别 $c_i = C_p(\mathbf{x}_i)$，而学习或训练过程 $L(\cdot)$ 的作用是确定分类器的参数 $p$。根据训练时是否提供、以及如何提供"正确答案"，学习方式可分为三类：监督学习（Supervised learning）依据每个样本给定的期望类别 $c_i^*$ 对学到的参数 $p$ 进行修正；强化学习（Reinforcement learning）只依据一个二值的奖惩信号 $r \in \{0,1\}$ 来修正参数，而不直接给出正确类别；无监督学习（Unsupervised learning）则完全不依赖任何期望类别的修正信号。

在具体方法层面，常用的分类器包括基于概率密度建模的贝叶斯分类、基于超平面划分的线性判别函数、基于一系列问答式规则的决策树，以及不依赖类别标签的聚类分析，此外还有大量其他方法。以下按统计方法、线性方法、非度量方法与无监督方法四条线索展开。

\subsection{统计方法}

统计方法的目标是做出统计意义上最优的决策：已知带标签数据集 $D$（分属若干类别）、类别数目 $c$ 以及问题的统计结构，输出对测试样本在统计意义上最优的分类结果。贝叶斯分类与最近邻规则是这一类方法中最具代表性的两种，此外还有最大似然估计、隐马尔可夫模型、密度估计等。这类方法常见于语音识别、文本识别、垃圾邮件过滤、目标检测等应用中。

以木材分类为例，令离散随机变量 $\omega$ 表示当前的真实状态：$\omega = \omega_1$ 表示白蜡木，$\omega = \omega_2$ 表示桦木。状态 $\omega_i$ 出现的先验概率记为 $P(\omega_i)$，满足 $P(\omega_i) \in [0,1]$ 且 $P(\omega_1) + P(\omega_2) = 1$。如果手头只有先验知识而没有任何观测，合理的决策是：当 $P(\omega_1) > P(\omega_2)$ 时判为 $\omega_1$，否则判为 $\omega_2$。

但通常还能观测到一个连续随机变量，例如亮度 $x$，其在给定类别下的分布由类条件概率密度 $p(x|\omega_i)$ 描述，表示在木材类型为 $\omega_i$ 的条件下观测到亮度 $x$ 的可能性。

若先验概率 $P(\omega_i)$ 与类条件概率 $p(x|\omega_i)$ 均已知，当观测到具体的 $x$ 时应如何决策？答案由贝叶斯规则给出：
\begin{equation}
    P(\omega_i|x) = \frac{p(x|\omega_i)P(\omega_i)}{p(x)}, \qquad p(x) = \sum_{j=1}^{2} p(x|\omega_j)P(\omega_j)
\end{equation}

这一规则把先验概率转化为观测到 $x$ 之后的后验概率 $P(\omega_i|x)$。据此可以给出朴素贝叶斯决策规则：若 $P(\omega_1|x) > P(\omega_2|x)$ 则判为 $\omega_1$，否则判为 $\omega_2$。这套方法的基本前提是所有相关概率均为已知，它本质上是针对两类特征分布存在重叠这一情形的统计处理方式，并且可以证明贝叶斯分类器能够使错误率最小化。

另一种不依赖概率模型、只依赖已标注样本本身的方法是最近邻规则。给定 $n$ 个带标签样本组成的集合 $D^n = \{\mathbf{x}_1, \mathbf{x}_2, \ldots, \mathbf{x}_n\}$，样本分属类别 $\omega_i$（$i = 1, \ldots, c$），对一个新的查询样本 $\mathbf{x}$，最近邻分类器在 $D^n$ 中找到与 $\mathbf{x}$ 距离最近的样本 $\mathbf{x}_i$，直接输出该样本的类别。这一规则在特征空间中天然对应一种沃罗诺伊镶嵌（Voronoi tessellation）：空间被划分为若干区域，每个区域内的任意一点都以某个训练样本为最近邻。将该思想推广，$k$ 近邻分类器在 $D^n$ 中找出与 $\mathbf{x}$ 最近的 $k$ 个样本（$k \le n$），选取这 $k$ 个样本中占多数的类别作为判决结果，从而在一定程度上抑制了单个近邻样本可能带来的噪声影响。

\subsection{线性方法}

线性方法的出发点是用一个判别函数 $g(\mathbf{x})$ 来刻画类别之间的边界，相比统计方法，它对数据分布的先验假设更少，几何直觉也更为直观：输入是带标签数据集 $D$ 与类别数 $c$，输出是一个判别函数 $g(\mathbf{x})$，问题的关键在于如何求得这个函数。常见方法包括感知机、Fisher 线性判别以及支持向量机，应用场景涵盖文档分类等高维分类问题。

对两个类别 $\omega_1, \omega_2$ 而言，一个在 $\mathbf{x}$ 各分量上呈线性形式的判别函数写作
\begin{equation}
    g(\mathbf{x}) = w_1 x_1 + \cdots + w_d x_d + w_0 = \mathbf{w}^t \mathbf{x} + w_0
\end{equation}

其中 $\mathbf{w} \ne 0$ 是权重向量，$w_0$ 是阈值。决策规则是：若 $g(\mathbf{x}) \ge 0$ 则判为 $\omega_1$，否则判为 $\omega_2$。方程 $g(\mathbf{x}) = 0$ 定义的超平面 $H$ 把特征空间划分成两个半空间，分别对应两个类别；权重向量 $\mathbf{w}$ 正是这个分界超平面的法向量，方向指向 $\omega_1$ 所在的决策区域。$g(\mathbf{x})$ 的取值还具有几何意义，等于样本点 $\mathbf{x}$ 到分界超平面的代数距离乘以 $\|\mathbf{w}\|$，即 $r = g(\mathbf{x}) / \|\mathbf{w}\|$。这种判别函数的优势在于：在特征分布满足一定条件时它是最优分类器，训练与计算都很简单，结构固定，也便于用专用硬件实现。

当两类样本在原始特征空间中并非线性可分时，支持向量机给出了一种处理办法：通过非线性核函数 $\varphi(\cdot)$ 把样本 $\mathbf{x}_k$ 映射到一个维度远高于原空间的新空间，记 $\mathbf{y}_k = \varphi(\mathbf{x}_k)$。理论上，只要映射选得恰当，两类样本在高维空间中总能被一个超平面分开。将阈值并入权重向量，记 $\mathbf{a} = (w_0\ \mathbf{w})^t$，并令 $y_0 = 1$，判别函数写作 $g(\mathbf{y}) = \mathbf{a}^t \mathbf{y}$。若用 $z_k = \pm 1$ 标记样本 $\mathbf{x}_k$ 属于 $\omega_1$ 还是 $\omega_2$，一个能够正确分开两类样本的超平面需要对所有 $k = 1, \ldots, n$ 满足 $z_k g(\mathbf{y}_k) \ge 1$。支持向量机的求解思路是调整 $\mathbf{a}$，使得间隔 $b$ 最大化，即让
\begin{equation}
    \frac{z_k g(\mathbf{y}_k)}{\|\mathbf{a}\|} \ge b, \quad \forall k
\end{equation}

间隔 $b$ 越大，分类器对未见样本的泛化能力通常也越好；真正决定分界超平面位置的，只是那些恰好落在间隔边界上的少数样本，即所谓的支持向量。核函数的选取、具体的训练算法以及样本非线性可分时的处理方式，是这一方法进一步需要考虑的问题。

\subsection{非度量方法}

当样本的属性是名义型（nominal）的、彼此之间既无法定义有意义的距离、也不存在天然的次序关系时，前述依赖距离或概率密度的方法便不再适用，这类问题需要用非度量方法处理。以水果分类为例，一个样本可能表示为若干离散属性的组合，例如西瓜是（颜色＝绿；大小＝大），西柚是（颜色＝黄；形状＝圆；大小＝大），柠檬是（颜色＝黄；形状＝圆；大小＝小）。此类方法在自动电话应答系统、定向广告预测等场景中有广泛应用，常见方法包括决策树、随机森林、字符串匹配以及基于规则的方法。

决策树以一系列问答的形式组织决策过程：内部节点代表针对某个属性的提问，例如"颜色？"；从节点引出的边代表该属性可能的取值，例如"绿、黄、红"；叶节点代表最终的类别，例如苹果、葡萄、樱桃。沿着从根节点到叶节点的路径依次回答问题，最终即可确定一个样本的类别。

构造决策树的输入是一个已预先分类好的训练集 $D$。若集合 $S \subseteq D$ 中所有样本都属于同一类别，则称 $S$ 是"纯"的（pure）。构造目标是得到一棵所有叶节点都纯的决策树，构造过程可归纳为一个递归算法：以整个训练集 $D$ 作为根节点开始，若当前节点上的样本集合是纯的，则该节点直接成为对应类别的叶节点；否则就要找出一个或多个属性，把当前集合划分到若干子节点中，并对每个子节点重复上述过程。

在构造决策树时需要考虑几个方面。其一是分支数目，即每个节点下设的子树个数，称为分支因子 $b$。其二是每次决策依赖的属性个数：只依据单一属性做决策称为单属性决策（monothetic），依据多个属性组合做决策则称为多属性决策（polythetic）。其三是节点的不纯度：节点 $N$ 的不纯度 $i(N)$ 在该节点内所有样本属于同一类别时为零，在各类别分布均匀时达到最大。选择用哪个属性划分节点时，若属性 $s$ 把节点 $N$ 划分为两个子节点 $N_L, N_R$，对应的样本比例为 $P_L, P_R$（满足 $P_L + P_R = 1$），则该次划分带来的不纯度下降量定义为
\begin{equation}
    \Delta i(s, N) := i(N) - \big(P_L\, i(N_L) + P_R\, i(N_R)\big)
\end{equation}

划分属性通常按照能最大化这一下降量的方向来选取。此外，何时停止继续划分、是否需要对已生成的树进行剪枝，也是构造过程中需要处理的问题。

只有一层、直接以一次属性判断划分出两个叶节点的决策树被称为决策树桩（decision stump），它是一种"弱学习器"，计算速度快、彼此相互独立，因此常被用作集成方法中的基础分类器。基于决策树桩等弱学习器，进一步衍生出两类重要的集成方法。随机森林由多棵互不相关的决策树组成：先通过有放回抽样（bagging）从原始训练集 $D$ 中抽出若干个规模更小的子训练集 $D' \subset D$（每个子集样本数 $n' < n$），再用每个子集分别训练一棵决策树作为组件分类器，最终对多棵树的判决结果取多数票，这种做法能够显著提升分类结果的稳定性。提升方法（boosting）则采用不同的思路：迭代地训练一系列弱分类器，每一轮都重点关注前一轮被错分的样本，AdaBoost 是这一思路的代表性算法。

\subsection{无监督学习}

与前述方法不同，无监督学习面对的是一组完全没有标签的训练样本，输入是 $n$ 个未标注样本，输出是一个能够识别测试样本的分类器（或者说，一种对数据内在结构的划分）。这类方法可以用来降低人工标注样本的成本、处理随时间变化的样本分类问题、为监督学习挑选合适的特征，也可以用来分析样本本身的性质、辅助选择恰当的分类方法。无监督学习的成功程度很大程度上取决于对问题所作的假设：数据分布的模型形式已知但参数未知、类别数目已知但分布未知、抑或类别数目和分布都完全未知，这几种情形对应着不同难度的问题变体。

$k$-均值聚类是最基本的无监督方法之一：输入是 $n$ 个未标注样本 $\mathbf{x}_1, \ldots, \mathbf{x}_n$ 以及类别数目 $c$（即 $k$），输出是 $c$ 个类别各自的均值向量 $\boldsymbol{\mu}_1, \ldots, \boldsymbol{\mu}_c$。算法的核心是交替更新每个样本的标签与每个类别的均值向量：先从样本中随机选取 $c$ 个初始均值，然后反复执行"依据与各均值的距离对所有样本重新分类"与"依据新的分类结果重新计算各均值"这两个步骤，直至均值不再发生变化为止。由于初始均值的选取带有随机性，同一批数据从不同初始值出发，迭代收敛的路径与最终结果也可能不同。

层次聚类给出了另一种组织样本的方式：若样本 $\mathbf{x}, \mathbf{x}'$ 在某一层级被归入同一簇，那么在更高的层级中它们也仍然属于同一簇。在层级 $k=1$ 时，$n$ 个样本各自独立成簇，共有 $c=n$ 个簇；层级 $k>1$ 时共有 $c = n-k+1$ 个簇；到层级 $k=n$ 时，全部样本合并为唯一一个簇。这种嵌套结构可以用树状图（dendrogram）直观呈现，树的高度对应某种相似度或距离度量，这既提供了衡量簇间相似程度的手段，也为估计数据中"自然"的簇数目提供了线索。

层次聚类的计算方向可以自顶向下（分裂式），也可以自底向上（凝聚式）：分裂式适合目标簇数远小于样本数的情形，凝聚式方法则相对简单，其思路是每一步都把当前距离最近的两个簇合并为一个新簇，如此反复，直至相似度出现较大幅度的下降为止。判断两个簇 $D_i, D_j$（分别含 $n_i, n_j$ 个样本，均值为 $\mathbf{m}_i, \mathbf{m}_j$）之间距离的方式并不唯一，例如最近邻聚类采用两簇中最近的一对样本点之间的距离
\begin{equation}
    d_{\min}(D_i, D_j) = \min_{\mathbf{x} \in D_i,\ \mathbf{x}' \in D_j} \|\mathbf{x} - \mathbf{x}'\|
\end{equation}

而最远邻聚类则采用两簇中最远的一对样本点之间的距离
\begin{equation}
    d_{\max}(D_i, D_j) = \max_{\mathbf{x} \in D_i,\ \mathbf{x}' \in D_j} \|\mathbf{x} - \mathbf{x}'\|
\end{equation}

不同的距离定义会导致凝聚过程形成不同形态的簇，这也是层次聚类在实际应用中需要根据数据特点加以选择的地方。